%% 10-intro.tex — Introduction

\section{Introduction}
\label{sec:introduction}

The Deutsche Digitale Bibliothek (DDB)\footnote{\url{https://www.deutsche-digitale-bibliothek.de}}
aggregates over 65 million cultural heritage objects from German archives,
libraries, museums, audiovisual archives, monument conservation offices, and
research institutions --- seven institutional sectors --- all encoded using the
DDB profile of the Europeana Data Model
(DDB-EDM)~\cite{ruehle2014ddbedm,edm2017definition}.
EDM achieves structural interoperability: every object, regardless of origin,
is typed as an \texttt{edm:ProvidedCHO} with a common property set drawn from
Dublin Core, DCMI Metadata Terms, and EDM extensions.
The cost of this uniformity is semantic flattening: the same property name
carries divergent semantics across sectors.

The \texttt{dc:type} property exemplifies the problem.
In the library sector, it holds terms from controlled bibliographic
vocabularies (e.g.\ \textit{Buch}, \textit{Monografie}); in the museum
sector, free-text descriptions of object type (e.g.\ \textit{Fotografie},
\textit{Zeichnung}); and in the archive sector, structural hierarchy
identifiers such as \textit{Bestand} or \textit{Archivale}.
A schema-level ontology alignment tool --- one that operates on lexical and
structural similarity between ontology class and property definitions --- sees
a single property \texttt{dc:type} and produces a single mapping.
It cannot recover the three distinct semantic regimes from the ontology
specification alone.
Corpus-based analysis, examining the actual value distribution of each
property stratified by sector and mediatype, is therefore not an optional
refinement but a precondition for reliable automated alignment.

This paper presents a corpus-driven three-layer methodology for aligning
DDB-EDM records to domain ontologies, mediated through
\textsc{mocho}~\cite{tan2026mocho}, a mid-level hub ontology that bridges EDM
to RDA~\todo{cite RDA}, RiC-O~\todo{cite RiC-O}, VRA~Core, and Music
Ontology~\cite{ceriani2018audio} via WEMI-level typing.
Three contributions are reported:

\begin{enumerate}
  \item \textbf{Corpus-driven methodology.}
        A three-layer dispatch architecture that profiles field distributions
        before constructing any alignment rule, and encodes alignment decisions
        as versioned lookup tables rather than hardcoded logic
        (Section~\ref{sec:methodology}).

  \item \textbf{Field profile and alignment table.}
        A data-derived field profile of 115{,}432 DDB records covering all
        seven sectors, and a corresponding alignment table mapping every
        observed (entity type, JSON key) pair to its target RDA property,
        WEMI level, and match method --- with gaps documented explicitly
        (Section~\ref{sec:results}).

  \item \textbf{Open-source ingest scripts.}
        A reproducible Python pipeline (\texttt{transform\_edm\_to\_mocho.py})
        that converts DDB JSON-LD records from the DDB API to N-Triples via
        the three-layer dispatch, together with all dispatch tables
        (\texttt{alignment\_ddbedm\_mocho.csv},
        \texttt{lookup\_htype\_doco\_rico.csv},
        \texttt{lookup\_dctype\_to\_class.csv}) \todo{confirm GitHub URL and licence}.
\end{enumerate}

\paragraph{Paper structure.}
Section~\ref{sec:related} reviews related alignment work.
Section~\ref{sec:methodology} describes the three-layer dispatch architecture.
Section~\ref{sec:results} reports coverage statistics on the pilot corpus.
Section~\ref{sec:discussion} discusses gaps and generalizability.
Section~\ref{sec:conclusion} concludes.
